\section{Relational Structures}\label{appendix:RS}

\body{
    A \define{signature} $\sigma$ is a set $S$ of relation symbols and an \define{arity} function $\text{ar}:S\rightarrow \mathbb{N}_0$.
    A (relational) \define{structure} $A$ over a signature $\sigma=(S,\text{ar})$ is a set $A$ along with an \define{interpretation function} $I$ which assigns to each symbol $R\in S$ a relation $R^A$ over $A$ of arity $\text{ar}(R)$. When the context is clear, we abuse notation and write $R$ for the relation $R^A$.
}

\begin{example}
    A directed graph can be encoded as a relational structure in several ways. For a simple directed graph $G$ with vertex set $V(G)$, we may use the signature $\sigma_{\mathrm{E}}=\{\mathsf{E}\}$, where $\operatorname{ar}(\mathsf{E})=2$, and interpret it on $V(G)$ by $\mathsf{E}^G=\{(u,v)\in V(G)^2 : \text{there is an edge from $u$ to $v$ in $G$}\}$. This encoding records adjacency but does not distinguish parallel edges.
\end{example}

\begin{example}
    Alternatively, a directed graph $G=(V,E,s,t)$ with loops and parallel edges may be encoded over the signature $\sigma_{\mathrm{inc}}=\{\mathsf{Vertex},\mathsf{Edge},\mathsf{Source},\mathsf{Target}\}$, where $\mathsf{Vertex}$ and $\mathsf{Edge}$ are unary and $\mathsf{Source}$ and $\mathsf{Target}$ are binary. The underlying set is the coproduct (i.e. disjoint union) $V+E$, with $\mathsf{Vertex}^G=V$, $\mathsf{Edge}^G=E$, $\mathsf{Source}^G=\{(e,s(e)):e\in E\}$, and $\mathsf{Target}^G=\{(e,t(e)):e\in E\}$. This encoding retains the identity of each edge and therefore distinguishes parallel edges.
\end{example}

%\begin{example} As another example, a simple university database may be regarded as a structure over the signature $\sigma_{\mathrm{univ}}=\{\mathsf{Student},\mathsf{Course},\mathsf{Enrolled}\}$, where the first two relation symbols are unary and $\mathsf{Enrolled}$ is binary. Its underlying set contains both students and courses, and $\mathsf{Enrolled}(x,y)$ records that the student $x$ is enrolled in the course $y$. \end{example}

\body{
    Given two structures $A$ and $B$ over a signature $\sigma=(S,\text{ar})$, a \define{homomorphism} from $A$ to $B$ is a function $f:A\rightarrow B$ so that for each relation $R\in S$ and each tuple $\vec{x}$ from $A$ with compatible arity, we have $R^A(\vec{x})\implies R^B(f(\vec{x}))$. Here $f(\vec{x})$ is computed pointwise. Relational structures and their homomorphisms form a category (c.f. Proposition~\ref{prop: relational structures category}). This category has all limits~\cite[Paragraph~3.10]{parker2022exponentiability}, which is enough for our purposes here.
}

\begin{proposition}\label{prop: relational structures category}
    The structures over a given signature $\sigma$, along with the corresponding homomorphisms between them, form a category with all limits.
\end{proposition}
\begin{proof}
    It is evident it forms a category: homomorphisms of structures consist of certain kinds of functions on the ground sets, so, once one realizes that relations are preserved under composition, associativity and identities are immediate. One can consult Parker~\cite[Paragraph~3.10]{parker2022exponentiability} to see it has limits; for completeness we will give details here. Let $D$ be a small diagram of $\sigma$-structures. Working in $\cat{Set}$, take the limit $L$ of the underlying point sets with projections $\pi_j\colon L\to\pts{D(j)}$. For each relation symbol $R\in S$ of arity $k$, define
    \[
        R^L=\{\vec{x}\in L^k:R^{D(j)}(\pi_j(\vec{x}))\text{ for every }j\}.
    \]
    By construction, the projections $\pi_j:L\rightarrow D(j)$ are homomorphisms and hence form a cone over $D$. Let $(X,\{q_j\})$ be any other cone over $D$. The universal property of the limit of the underlying point sets gives a unique function $q\colon\pts{X}\rightarrow L$ such that $\pi_j\circ q=q_j$ for every $j$. If $R^X(\vec{x})$, then $R^{D(j)}(q_j(\vec{x}))$ for every $j$, since each $q_j$ is a homomorphism. It follows from the definition of $R^L$ that $R^L(q(\vec{x}))$, and hence $q$ is a homomorphism. Therefore $L$ is the limit of $D$.
\end{proof}

\subsection{Modeling Equality}
\label{subsec:model_eq}

\body{In the framework of relational structures, we cannot ask directly about equality of points in a first-order formula. However, our methods extend in the following way: We can model equality using an additional binary relation that appears as a loop on every point in the input co-decomposition/structure. We note that it is possible for there to be structures in our calculations (in particular inside the heads of hypotheses), for which this binary relation relates two distinct points. This technique is made precise in Observation~\ref{loop_trick}.}

\body{For a structure $G$, let $G^\circ$ be the structure obtained from $G$ by adding a new relation $\circ\subseteq\pts{G}^2$ which is reflexive on every point but relates no distinct points. That is, it is a loop everywhere. For a co-decomposition $\Ta$, we denote by $\Ta^\circ$ the co-decomposition obtained from $\Ta$ by replacing each structure by its variant with this new relation.}

\begin{observation}\label{loop_trick}
   Let $\Ta$ be a co-decomposition. Then $\lim(\Ta^\circ)=(\lim(\Ta))^\circ$. 
\end{observation}

\body{We can use our machinery on $\Ta^\circ$ to ask questions about equality on $\lim\Ta$. That is, evaluating FO with equality requires only a linear amount of pre-processing on the co-decomposition.}